Literature Gap Summary for Research Design: ESG ABSA with Tone-Aware, Multilingual, Explainable Pipelines

Executive framing
Across the reviewed literature on ESG ABSA, a coherent, end-to-end, bilingual (Indonesian/English) pipeline that delivers sentence- and region-grounded ESG evidence with explicit tone distinctions (commitment, action, outcome) and regulator-aligned outputs remains elusive. The common thread is a fragmentation of approaches: many studies address either sentence-level ABSA, or cross-lingual ABSA, or climate-domain ABSA, or ethics and explainability in isolation, but few integrate all these threads into a single auditable architecture suitable for OJK/IDX and international investors. Moreover, although numerous domain-adapted models (ClimateBERT, FinBERT, ESG-BERT) show promise, there is limited evidence on how to orchestrate multiple LLMs and prompt templates in a controlled, reproducible manner for ESG disclosures that include multimodal elements (tables, figures) and non-English language, while guaranteeing explainability and regulatory traceability. This synthesis identifies how these gaps cohere and why a dedicated research design is necessary to advance robust, trustworthy ESG ABSA in practice, and it motivates the subsequent methodology chapter.

1) Gap: Lack of an integrated, end-to-end, bilingual ESG ABSA pipeline with tone awareness

What the literature shows
Numerous works demonstrate ABSA at the sentence level or across languages, or specialize in climate/disclosure domains, yet few deliver a single, unified pipeline that (i) processes bilingual Indonesian/English sustainability reports, (ii) grounds analysis in a formal tone taxonomy (commitment, action, outcome), (iii) aligns extractions to a joint ESG ontology anchored to GRI/SASB/TCFD, and (iv) produces regulator-ready, auditable outputs end-to-end. Multiple strands exist in isolation (monolingual ABSA, cross-lingual ABSA, climate-target ABSA, and ontology-based extraction), but there is sparse evidence of seamless integration with end-to-end evaluation on Indonesian disclosures (Seow, 2025; , (Šmíd & Král, 2025; , Gorovaia & Makrominas, 2024; , Hang, 2025; , Heever et al., 2024; , Garrido-Merchán et al., 2023; , Hassani et al., 2025; , Schimanski et al., 2023), Kim et al., 2024).
What’s missing
A cohesive architecture that transforms PDFs or scans into sentence- and region-grounded ESG records with provenance, performs ABSA with a four-dimensional signal (aspectterm, sentimentlabel, tonetype, regulatoryanchor), and links outputs to GRI/SASB/TCFD concepts within a bilingual ontology. There is also a lack of end-to-end evaluation on Indonesian disclosures under OJK/IDX with bilingual performance metrics and regulatory traceability.
Implications
Without an integrated pipeline, comparisons across firms and jurisdictions are unreliable, auditability is weak, and the ability to satisfy investor due diligence and regulator inquiries is compromised. This gap motivates a design that combines OCR, multilingual ABSA, ontology grounding, and explainable outputs in a reproducible framework.
2) Gap: Insufficient emphasis on a formal tone taxonomy (commitment/action/outcome) as a diagnostic signal for credibility and greenwashing

What the literature shows
Most ABSA studies treat sentiment as a single polarity per aspect, with limited attention to epistemic modality or the regulatory relevance of forward-looking commitments versus realized outcomes. Although some climate-domain work investigates targets (net-zero) and commitments, there is no widely adopted, validated tone taxonomy that explicitly discriminates commitment, action, and outcome across ESG facets in a bilingual regulatory context (Seow, 2025; , Gorovaia & Makrominas, 2024; , Heever et al., 2024; , Garrido-Merchán et al., 2023; , Hassani et al., 2025; , Schimanski et al., 2023), Kim et al., 2024).
What’s missing
A validated tone taxonomy that (i) segments statements by tone type (commitment, action, outcome) and (ii) ties tone to regulatory anchors and measurable KPIs across GRI/SASB/TCFD concepts, with bilingual annotation guidelines and robust inter-annotator agreement metrics. There is also a need for corpus resources annotated with tone types in both Indonesian and English ESG disclosures.
Implications
The absence of a formal tone taxonomy handicaps the detection of greenwashing and weakens auditability. Embedding a tone-aware ABSA framework enables more precise credibility assessment and supports governance signaling research aligned with stakeholder theory and legitimacy theory.
3) Gap: Insufficient cross-language, regulator-aligned ontology integration for Indonesian/English ESG disclosures

What the literature shows
Cross-lingual ABSA literature demonstrates the feasibility of multilingual pipelines, but Indonesian-specific regulatory lexicons and bridging Indonesian-English ESG vocabulary to GRI/SASB/TCFD nodes remain underdeveloped. Ontology-based extraction in ESG is gaining traction, yet many efforts rely on surface lexicons rather than formal, machine-interpretable ontologies with cross-language mappings and provenance. There is a paucity of bilingual corpora annotated to ontology nodes and lack of standardized mappings to multiple frameworks (GRI, SASB, TCFD) in Indonesian contexts (Šmíd & Král, 2025; , Gorovaia & Makrominas, 2024; , Hang, 2025; , Kansal et al., 2024; , JIA et al., 2025; , Brauwers & Frăsincar, 2022).
What’s missing
A bilingual ESG ontology with Indonesian-English term mappings, cross-language synonym sets, and formal alignment to GRI/SASB/TCFD indicators. A pipeline that uses this ontology to drive ABSA, ensure cross-language consistency, and maintain regulatory traceability across Indonesian and international standards.
Implications
Without a robust bilingual ontology, multilingual ABSA will struggle with term drift, inconsistent aspect labeling, and misalignment with regulatory requirements. The ontology is essential for explainability, cross-document comparability, and regulator-friendly outputs.
4) Gap: Underdeveloped multimodal and document-structure aware ESG ABSA (OCR, tables, charts, layout)

What the literature shows
ESG reports are multimodal, containing narratives, tables with KPI data, and charts. While multimodal NLP and document understanding are advancing, there is still a lack of end-to-end ESG ABSA pipelines that integrate OCR quality, table/chart extraction, and cross-language grounding within an ontology-driven ABSA framework. The literature highlights the risk that OCR errors propagate into ABSA and compromise sentiment and regulation mappings, especially in long, complex sustainability documents (Gao, 2025; , Wang, 2025; , Gu et al., 2022; , Aggarwal et al., 2023; , Li et al., 2024; , Wang et al., 2023; , Yu, 2025).
What’s missing
An integrated pipeline that preserves document structure via page-batching, performs layout-aware OCR, extracts and aligns tables/charts with textual sentiment, and jointly reasons across modalities to produce sentence-level ABSA with regulatory tags. Evaluation protocols should quantify how OCR and multimodal extraction quality affect ABSA outcomes and how cross-language content (Indonesian/English) is handled in tables and figures.
Implications
Without multimodal integration, ABSA will miss or misinterpret essential evidentiary signals, undermining regulatory alignment and stakeholder trust. A multimodal ABSA design is critical for credible ESG analysis in Indonesian markets.
5) Gap: Limited reproducibility, explainability, and regulator-facing outputs in ESG ABSA with ground-truth labeling scarcity

What the literature shows
While XAI techniques (LIME, SHAP, attention visualization) are increasingly applied in ESG contexts, many studies do not publish complete reproducibility artifacts (prompts, schemas, ground-truth corpora) or provide regulator-oriented explanations that join textual signals with regulatory mapping. The absence of expert-labeled ground truth in many ESG ABSA settings constrains rigorous evaluation and replication, making it difficult to compare methods or demonstrate auditability to regulators and investors (Seow, 2025; , Moreno & Caminero, 2020; , Srivastava & Anand, 2023; , Lin et al., 2024; , Cardoni et al., 2019; , Yang, 2025).
What’s missing
A comprehensive reproducibility package including seven prompt templates, two baseline models, an explicit JSON schema for outputs, bilingual annotated corpora with tone labels and ontology mappings, and publicly accessible evaluation protocols. A clear plan for expert-in-the-loop labeling and incremental ground-truth generation over time is needed.
Implications
In the absence of ground-truth labels, the design must rely on robust diagnostics, disagreement analyses, and cross-model validation to triangulate performance, while building a scalable plan to generate expert-labeled data that will eventually anchor evaluation.
6) Gap: Experimental design and evaluation paradigms without expert-labeled ground truth

What the literature shows
A number of ESG ABSA studies employ supervised learning on domain-specific labels, yet many lack expert-labeled ground truth that would serve as a gold standard for sentence-level tone and ontology mappings. This makes it difficult to quantify absolute performance and to validate models for regulatory credibility. There is also a need for cross-language evaluation metrics and for domain-aware calibration of model confidence in bilingual settings (Seow, 2025; , Gorovaia & Makrominas, 2024; , Hang, 2025; , Heever et al., 2024; , Garrido-Merchán et al., 2023; , Hassani et al., 2025; , Schimanski et al., 2023).
What’s missing
An experimental design that (i) acknowledges the absence of a full ground-truth, (ii) uses robust proxy metrics (disagreement analyses, calibration curves, cross-language concordance), (iii) includes a plan for expert-driven ground-truth annotation in stages, and (iv) defines a reproducibility plan that ensures future analysts can rebuild the experiments.
Implications
The research must embrace a flexible, exploratory design that can generate interim, task-specific ground-truth data while providing rigorous diagnostics for reliability, stability, and interpretability.
7) How this thesis positions to fill the gaps

Integrated, explainable, tone-aware bilingual ABSA: This research proposes a holistic pipeline that ingests PDFs or scans, preserves document structure via page-batching, performs sentence- and region-level ABSA with a four-field signal (aspectterm, sentimentlabel, tonetype, regulatoryanchor), and grounds outputs to a bilingual ESG ontology aligned to GRI/SASB/TCFD. It introduces a formal tone taxonomy (commitment, action, outcome) as a core methodological contribution for automated greenwashing analysis, and utilizes a knowledge-graph grounding layer to support explainability and regulator-facing outputs.
Multi-model and prompt-template comparative strategy: The design explicitly incorporates multiple LLMs (e.g., generic multilingual encoders and climate/domain-adapted models) and seven prompt templates to explore variance, enable ensemble strategies, and build a robust diagnostic framework. This comparative strategy is the primary mechanism to cope with the absence of expert-labeled ground truth by triangulating signals across models and prompts.
Ground-truth development plan: A staged approach is proposed to generate expert-labeled data over time, starting with bilingual sentence- and clause-level annotations for a curated corpus of Indonesian/English ESG disclosures, focusing on tone labels and ontology mappings. This plan acknowledges the current ground-truth gap and provides a concrete trajectory toward expert-validated evaluation.
Reproducibility and regulator alignment: The design emphasizes reproducibility artifacts (PROMPTS, JSON schemas, model configurations, prompts/templates, evaluation scripts) and regulator-friendly explainability (ontology paths, provenance tracking, finance/regulatory mappings), ensuring outputs are auditable and actionable for investors and regulators.
Transition to methodology
The following methodology chapter will detail data collection, annotation schemas, the end-to-end pipeline architecture, the seven prompt templates, model configurations, ontology mapping strategies, evaluation plans (including disagreement analyses and calibration metrics), and the reproducibility package. It will also outline the ground-truth generation plan and staged expert-in-the-loop labeling protocol to move from exploratory evaluation toward a formal gold standard.
References and grounding

Ground the design in ABSA literature across domains (environmental, financial, climate), cross-language ABSA, and ontology-driven information extraction; incorporate XAI approaches (LIME, SHAP, attention visualization) for regulator-facing interpretability; and anchor to governance theories (stakeholder theory, legitimacy theory) to justify tone-focused analysis and greenwashing detection. The seven prompt templates and two model configurations will be grounded in the cited works on prompt variation, RAG-augmented extraction, and domain-adapted models (ClimateBERT, FinBERT, ESG-BERT). Where possible, explicit DOIs and venues should be inserted in the final manuscript.
Note: If you provide the exact seven prompt templates and the two models you plan to use, I can tailor the justification to your concrete setup and generate a polished, IEEE-formatted References section with precise DOIs and bibliographic details.